\heading{量子力学A-23秋-期中}
\noteinfo{整理自 Quantum Mechanics: Fall 2023 Midterm Exam: Brief Solutions；题面和解答按项目格式重写。}

\begin{question}
一维谐振子
\[
  H={p^2\over2m}+{m\omega^2x^2\over2}.
\]
初态为
\[
  \phi(x,0)=A\left(x+\sqrt{\hbar\over m\omega}\right)^2
  \exp\left(-{m\omega x^2\over2\hbar}\right).
\]
\begin{enumerate}
  \item 求归一化常数 $A$。
  \item 测量能量的可能结果及概率。
  \item 求 $\phi(x,t)$。
  \item 求 $\expval{x},\expval{p},\expval{x^2},\expval{p^2}$ 并检验不确定关系。
\end{enumerate}
\begin{solution}
记
\[
  x=\sqrt{\hbar\over2m\omega}(a+a^\dagger),\qquad
  \psi_0(x)=\left({m\omega\over\pi\hbar}\right)^{1/4}
  \exp\left(-{m\omega x^2\over2\hbar}\right).
\]
于是
\[
  \phi(x,0)=A{\hbar\over2m\omega}\left({\pi\hbar\over m\omega}\right)^{1/4}
  \left(\sqrt2\ket2+2\sqrt2\ket1+3\ket0\right).
\]
归一化给
\[
  |A|=\left({m\omega\over\pi\hbar}\right)^{1/4}{2m\omega\over\hbar\sqrt{19}}.
\]
取 $A$ 为正实数，则
\[
  \ket{\phi(0)}={1\over\sqrt{19}}\left(\sqrt2\ket2+2\sqrt2\ket1+3\ket0\right).
\]
测量能量的结果和概率为
\[
  {1\over2}\hbar\omega:\ {9\over19},\qquad
  {3\over2}\hbar\omega:\ {8\over19},\qquad
  {5\over2}\hbar\omega:\ {2\over19}.
\]
时间演化为
\[
  \ket{\phi(t)}={\ee^{-\ii\omega t/2}\over\sqrt{19}}
  \left(\sqrt2\,\ee^{-2\ii\omega t}\ket2
  +2\sqrt2\,\ee^{-\ii\omega t}\ket1+3\ket0\right).
\]
用升降算符直接计算得
\[
  \expval{x}={20\sqrt2\over19}\sqrt{\hbar\over2m\omega}\cos\omega t,
\]
\[
  \expval{p}=-{20\sqrt2\over19}\sqrt{\hbar m\omega\over2}\sin\omega t,
\]
\[
  \expval{x^2}={\hbar\over2m\omega}{43+12\cos2\omega t\over19},
  \qquad
  \expval{p^2}={\hbar m\omega\over2}{43-12\cos2\omega t\over19}.
\]
因此
\[
  \sigma_x^2\sigma_p^2=\left(\expval{x^2}-\expval{x}^2\right)
  \left(\expval{p^2}-\expval{p}^2\right)\ge{\hbar^2\over4}.
\]
\end{solution}
\end{question}

\begin{question}
考虑
\[
  H={p^2\over2m}+{\hbar^2\over2m}\beta_1[\delta(x-L)+\delta(x+L)]
  -{\hbar^2\over2m}\beta_2\delta(x),
\]
其中 $\beta_1,\beta_2,L>0$。
\begin{enumerate}
  \item 定性画出束缚态。
  \item 推导束缚态能量方程。
  \item 判断出现束缚态的参数条件。
\end{enumerate}
\begin{solution}
束缚态能量写成 $E=-\hbar^2\kappa^2/(2m)$，$\kappa>0$。势能为偶函数，态可取偶或奇。奇态在 $x=0$ 为零，中央吸引 $\delta$ 不作用；两侧 $\delta$ 势为排斥，因此不存在奇束缚态。偶束缚态可写成
\[
  \psi(x)=\begin{cases}
  A\ee^{\kappa |x|}+B\ee^{-\kappa |x|},& |x|<L,\\
  C\ee^{-\kappa(|x|-L)},& |x|>L.
  \end{cases}
\]
边界条件为
\[
  \psi'(+0)-\psi'(-0)=-\beta_2\psi(0),
  \qquad
  \psi'(\pm L^+)-\psi'(\pm L^-)=\beta_1\psi(\pm L).
\]
消去系数得
\[
  \ee^{2\kappa L}={\beta_1(\beta_2+2\kappa)\over(\beta_2-2\kappa)(\beta_1+2\kappa)}.
\]
等价地
\[
  2\kappa L=\log{1+2\kappa/\beta_2\over(1-2\kappa/\beta_2)(1+2\kappa/\beta_1)}.
\]
由 $\kappa\to0^+$ 的阈值条件可得存在束缚态当且仅当
\[
  L-{2\over\beta_2}+{1\over\beta_1}>0,
  \qquad\text{即}\qquad
  \beta_2>{2\beta_1\over1+\beta_1L}.
\]
当 $L\to0$ 时，这化为单个 $\delta$ 势的条件 $\beta_2>2\beta_1$。
\end{solution}
\end{question}

\begin{question}
氢原子电子。设轨道波函数
\[
  \psi(\vb r)=A(x+y+z)\ee^{-r/2a}.
\]
\begin{enumerate}
  \item 求归一化常数 $A$。
  \item 求 $L_x,L_y,L_z$ 的期望。
  \item 若电子自旋 $S=1/2$ 且总态为 $\psi(\vb r)\ket{\downarrow}$，定义 $\vb J=\vb L+\vb S$。测量 $J^2,J_z$，求可能结果、概率和塌缩态。
\end{enumerate}
\begin{solution}
由球面对称性，
\[
  \int (x+y+z)^2\ee^{-r/a}\dd^3r
  =\int r^2\ee^{-r/a}\dd^3r=96\pi a^5,
\]
故
\[
  |A|={1\over\sqrt{96\pi}\,a^{5/2}}.
\]
波函数可取为实函数，且在 $x\to y\to z\to x$ 的循环转动下不变，所以
\[
  \expval{L_x}=\expval{L_y}=\expval{L_z}=0.
\]
轨道态展开为
\[
  \ket{\psi}={1\over\sqrt3}\left(-\ee^{-\ii\pi/4}\ket{1,1}
  +\ket{1,0}+\ee^{\ii\pi/4}\ket{1,-1}\right),
\]
其中 $\ket{1,m}$ 表示 $\psi_{21m}$。用 $\ell=1$ 与 $s=1/2$ 的 C-G 系数展开 $\ket{\psi}\ket{\downarrow}$，测量结果为
\[
\begin{array}{c|c|c}
(J^2,J_z)&P&\text{塌缩态}\\ \hline
\left({15\over4}\hbar^2,{1\over2}\hbar\right)&{1\over9}
&{1\over\sqrt3}\left(\ket{1,1}\ket{\downarrow}+\sqrt2\ket{1,0}\ket{\uparrow}\right)\\
\left({3\over4}\hbar^2,{1\over2}\hbar\right)&{2\over9}
&{1\over\sqrt3}\left(\sqrt2\ket{1,1}\ket{\downarrow}-\ket{1,0}\ket{\uparrow}\right)\\
\left({15\over4}\hbar^2,-{1\over2}\hbar\right)&{2\over9}
&{1\over\sqrt3}\left(\sqrt2\ket{1,0}\ket{\downarrow}+\ket{1,-1}\ket{\uparrow}\right)\\
\left({3\over4}\hbar^2,-{1\over2}\hbar\right)&{1\over9}
&{1\over\sqrt3}\left(\ket{1,0}\ket{\downarrow}-\sqrt2\ket{1,-1}\ket{\uparrow}\right)\\
\left({15\over4}\hbar^2,-{3\over2}\hbar\right)&{1\over3}
&\ket{1,-1}\ket{\downarrow}
\end{array}
\]
\end{solution}
\end{question}

\begin{question}
粒子处于谐振势与有限高势垒组合势中：$|x|<a$ 内近似为谐振势，外部为有限常数势垒。讨论低能束缚态。
\begin{enumerate}
  \item 定性画出基态、第一激发态、第二激发态波函数。
  \item 写出内外区级数解的递推关系与边界匹配条件。
  \item 说明与普通谐振子的递推关系有何联系。
\end{enumerate}
\begin{solution}
低能束缚态在阱内类似谐振子本征态，在阱外指数衰减；基态偶且无节点，第一激发态奇且一个节点，第二激发态偶且两个节点。势能为偶函数，解可按宇称分类。

阱内把波函数写作
\[
  \psi_{\rm in}(x)=\ee^{-m\omega x^2/(2\hbar)}\sum_j c_{{\rm in},j}x^j,
\]
阱外写作指数衰减乘级数。代入定态薛定谔方程得到与谐振子相同形式的三项递推，只是能量不再由级数截断决定，而由 $x=\pm a$ 处
\[
  \psi_{\rm in}=\psi_{\rm out},\qquad
  \psi_{\rm in}'=\psi_{\rm out}'
\]
共同决定。有限势垒使能级相对无限谐振子发生偏移，且波函数在外部有尾部。
\end{solution}
\end{question}
